\begin{musica}{Mulher}{(desconhecido)}{Tuna de Medicina da Universidade do Porto}
\nota{Introdução: Gm A D B7 Em A D A7 A+5}

Mu\ac{D}{lher}, \ac{F#}{}__ mal que por meu \ac{Bm}{mal} és o meu \ac{B7}{bem}
Qual\ac{Em}{quer} \ac{B7}{}__ sabe que és le\ac{Em}{al} só quando és mãe.
É \ac{Gm}{lou}co, é sandeu quem dis\ac{A}{ser}
À mu\ac{D}{lher} seja ela qual \ac{B7}{for}
Mu\ac{Em}{lher}, tu és o a\ac{A}{mor}. \ac{A+5}{}___

Não, não és, tu és só traição de lés a lés,
Mal vês preso um coração pisá-lo aos pés,
Em cem como tu conto as cem
Entre as falsas, pois, cego eu não sou,
Mas diz-me outra vez: “amor vem!”
Diz-me, que eu vou.

\nota{Instrumental: D F# Bm B7 Em B7 Em}
\nota{\phantom{Instrumental: }Gm A D B7 Em A D A7 A+5}

Não, não és, tu és só traição de lés a lés,
Mal vês preso um coração pisá-lo aos pés,
Em cem como tu conto as cem
Entre as falsas, pois, cego eu não sou,
Mas diz-me outra vez: “amor vem!”

Diz-me, que eu \ac{Gm}{vou}. \ac{A D Dmaj7}{}_ _ ____

\posicoes{D: xx0232, Dmaj7: xx0222, Em: 022000, F#: 244322, Gm: 355333, A: x02220,
  A7: x02020, A+5: x02221, B7: 021202, Bm: x24432}
\end{musica}
